% Graph-backed presentation of a least counterexample.
%
% Companion to foundations/map-euler-foundations.tex and
% foundations/tait-reduction.tex.  Status markers as in those files.

\section{Graph-backed presentation, and why minimality is not optional}
\label{sec:graph-backed}

The compositional machinery is stated over rotation systems; the classical
bridge of Section~\ref{sec:tait-reduction} is stated over simple graphs. To
join them one needs a bridgeless spherical cubic map to be presentable as
the rotation system of a \emph{simple} graph.

The tempting statement is that every such map is so presentable. That
statement is false, and the counterexample is small enough to be worth
stating first, because it tells us exactly which hypothesis is doing the
work.

\subsection{The obstruction}

\begin{example}[the spherical theta multigraph]\label{ex:theta}
\Sprose
Let $D = \{u_0, u_1, u_2, v_0, v_1, v_2\}$, let the edge flip exchange $u_i$
with $v_i$, and let the vertex rotations be
\[
  \sigma_u = (u_0\,u_1\,u_2), \qquad \sigma_v = (v_0\,v_2\,v_1).
\]
The face permutation is then $(u_0\,v_2)(u_1\,v_0)(u_2\,v_1)$, so
\[
  V = 2, \qquad E = 3, \qquad F = 3, \qquad V - E + F = 2 .
\]
This map is spherical, cubic, connected, and bridgeless. But no simple
graph on two vertices is cubic --- a simple graph on two vertices has
maximum degree one --- so there is no equal-vertex-count presentation of it
by a simple graph.
\end{example}

Two vertices joined by three parallel edges: the multiplicity is exactly
what a simple graph cannot express. The unrestricted presentation claim is
therefore false, and any development asserting it is asserting something
untrue.

The repair is not to weaken the conclusion but to remember the hypothesis
that was available all along. We do not need to present \emph{every} map. We
need to present a \emph{least counterexample}, and the theta map is not one:

\begin{remark}\label{rem:theta-colourable}
\Sprose
The theta map is Tait colourable. Its three edges are incident with both
vertices, so assigning them three distinct colours properly colours both
stars at once. Hence it is not a counterexample to anything, and no
statement restricted to counterexamples is threatened by it.
\end{remark}

That is the shape of the whole section: multiplicity is exactly what
minimality removes.

\subsection{Setting}

Fix a map $M$ that is connected, spherical, cubic, and bridgeless, and that
is a \emph{least counterexample}: $M$ admits no Tait colouring, while every
map in the same class with strictly fewer vertices does.

\begin{lemma}[no loops]\label{lem:no-loops}
\Sprose
$M$ has no loop. This uses bridgelessness only; minimality is not needed.
\end{lemma}

\begin{proof}
Suppose $v$ carries a loop $\ell$. The loop occupies two of the three darts
at $v$, so exactly one further edge $e$ is incident with $v$, and its other
end is some vertex $w$. If $w = v$ then $e$ is a second loop at $v$, giving
$v$ four darts and contradicting cubicity; so $w \neq v$ and, in
particular, $M$ has at least two vertices.

Delete $e$. The only remaining edge at $v$ is the loop $\ell$, so $\{v\}$ is
a connected component of $M - e$, while $w \neq v$ lies in another. Thus
deleting $e$ increased the number of components: $e$ is a bridge,
contradicting bridgelessness.
\end{proof}

\subsection{Removing parallel edges}

This is where minimality is spent.

\begin{lemma}[the digon splice]\label{lem:digon-splice}
\Sconditional{Theorem~\ref{thm:digon-combinatorial}}
Suppose distinct vertices $u, v$ of $M$ are joined by exactly two parallel
edges, and $M$ is not the theta map.
Let $a = uu'$ and $b = vv'$ be the remaining edges at $u$ and at $v$. Then
$u' \neq v'$, and the map $M'$ obtained by deleting $u$, $v$, the two
parallel edges, $a$ and $b$, and inserting a single new edge $c = u'v'$, is
again connected, spherical, cubic, and bridgeless, with exactly two fewer
vertices.
\end{lemma}

\begin{proof}
Each of $u$ and $v$ has three darts, two consumed by the parallel pair, so
$a$ and $b$ are well defined and unique. By Lemma~\ref{lem:no-loops},
$u' \neq u$ and $v' \neq v$.
Moreover $u'\neq v$ and $v'\neq u$: either equality would make the remaining
edge a third $uv$ edge; cubicity and connectedness would then make $M$ the
theta map, which has been excluded. Thus the splice really removes $u,v$ and
joins two surviving vertices.

\emph{First, $u' \neq v'$.} Suppose $u' = v' = w$. Then $w$ is incident with
$a$ and with $b$, and being cubic it has exactly one further edge $f$. Now
$\{u,v,w\}$ meets the rest of $M$ only along $f$: the parallel pair, $a$ and
$b$ are internal to it, and $u$, $v$ have no other darts. Deleting $f$
therefore separates $\{u,v,w\}$ from the remainder, so $f$ is a bridge
unless there is no remainder --- and if there is none, $w$'s third dart has
nowhere to go without violating cubicity. Either way we contradict a
standing hypothesis, so $u' \neq v'$.

\emph{Cubicity.} Every vertex other than $u',v'$ is untouched. Each of $u'$
and $v'$ loses one incident edge ($a$, resp.\ $b$) and gains one ($c$), so
retains degree three.

\emph{Connectedness.} Any path of $M$ through the deleted region entered and
left along $a$ and $b$, and now traverses $c$ instead.

\emph{Bridgelessness.} Suppose $c$ were a bridge of $M'$, separating $M'-c$
into a part $A \ni u'$ and a part $B \ni v'$. Then in $M$ the only
connections between $A$ and $B$ run through $u$ and $v$, hence through $a$
and the parallel pair and $b$. Deleting $a$ from $M$ would then separate
$A$ from the rest, making $a$ a bridge of $M$ and contradicting
bridgelessness. For an edge $e\neq c$ of $M'$, expand every traversal of $c$
to the path $u' a u,e_1,v b v'$ in $M$, and contract every traversal of the
deleted two-vertex network in $M$ back to $c$. These operations preserve
walks after deleting $e$. Hence, if $M'-e$ were disconnected, then $M-e$
would be disconnected as well, contrary to bridgelessness of $M$.

\emph{Sphericity.} The count is consistent: $V' = V-2$, $E' = E-3$, and one
face is destroyed --- the face bounded by the parallel pair --- so
$F' = F-1$ and
\[
  V' - E' + F' = (V-2)-(E-3)+(F-1) = V - E + F = 2 .
\]
The step that must be discharged is the claim that the parallel pair bounds
a face --- it does not follow from the counts --- and
Theorem~\ref{thm:digon-combinatorial} supplies it from bridgelessness, Euler's
inequality, and a finite dart computation. Thus the splice uses no separation
theorem.
\end{proof}

\subsection{What the digon face really needs}

The splice above rests on the parallel pair bounding a face. That claim
decomposes further than it first appears, and the decomposition is worth
recording because half of it is free.

Write the digon's darts as $p_1, p_2$ at $u$ (on the two parallel edges
$e_1, e_2$) and $q_1, q_2$ at $v$, with $\alpha p_i = q_i$, and let $p_a$,
$q_b$ be the third darts at $u$ and at $v$.

\begin{lemma}[a cubic digon always shares a face]\label{lem:digon-shares}
\Sverified{Mettapedia.GraphTheory.FourColor.GoertzelV24BridgeTwoSided.sameCycle_digon}
In a cubic map, some face orbit contains a dart of $e_1$ and a dart of
$e_2$. No sphericity and no counting are needed.
\end{lemma}

\begin{proof}
The rotation $\sigma_u$ is a three-cycle on $\{p_1, p_2, p_a\}$, so either
$\sigma p_1 = p_2$ or $\sigma p_1 = p_a$. In the first case
$\varphi(q_1) = \sigma(\alpha q_1) = \sigma p_1 = p_2$, so $q_1$ and $p_2$
lie on one face. In the second case $\sigma_u = (p_1\,p_a\,p_2)$, whence
$\sigma p_2 = p_1$ and $\varphi(q_2) = p_1$, so $q_2$ and $p_1$ lie on one
face. With only three darts at a cubic vertex, two of them are cyclically
adjacent one way or the other, and there is no third possibility. The
formal statement takes that dichotomy as its hypothesis, which is precisely
what a three-dart rotation supplies.
\end{proof}

\begin{remark}[where sphericity actually enters]\label{rem:digon-exact}
What the splice needs is stronger: that the shared face is \emph{exactly}
the digon, with boundary $\{e_1, e_2\}$ and nothing more, since only then
does deleting the pair destroy precisely one face and keep the Euler count
of Lemma~\ref{lem:digon-splice} honest. Continuing the computation above in
the first case, $\varphi(p_2) = \sigma(q_2)$, and this returns to $q_1$ --
closing the two-cycle -- exactly when $\sigma q_2 = q_1$, that is, when the
rotation at $v$ is oriented compatibly with the one at $u$. Both
orientations give legitimate rotation systems; they differ in genus. So
sphericity is precisely the hypothesis that rules the incompatible one out,
and this is the point at which the digon lemma stops being free.

A small exhaustive check corroborates this before any proof is attempted.
Take the cubic multigraph on $\{u,v,w,x\}$ with a parallel pair joining $u$
to $v$, a parallel pair joining $w$ to $x$, and single edges $uw$ and $vx$;
enumerate all $2^4 = 16$ rotation systems on it. Of these, four are spherical
and twelve have genus one. In all four spherical systems the pair $\{e_1,
e_2\}$ bounds a face of length exactly two; among the genus-one systems only
four do, and eight do not. So sphericity is sufficient here and visibly not
necessary --- which is the asymmetry the remark predicts, and the direction
the splice actually consumes. This is corroboration on one family rather than
a proof, and the lemma keeps its \textsc{open} marker accordingly.
\end{remark}

\subsection{The digon face, from bridgelessness}

The exhaustive check suggests the claim; here is the argument. It is worth
setting down because it consumes no new hypothesis: the third edge at $u$
would be a bridge, and bridgelessness is already standing.

\begin{lemma}[a parallel pair bounds a face]\label{lem:digon-face}
\Sconditional{sphere separation (Section~\ref{sec:tait-reduction})}
Let $M$ be connected, spherical, cubic and bridgeless, let distinct vertices
$u,v$ be joined by parallel edges $e_1,e_2$, and suppose $M$ is not the theta
map of Example~\ref{ex:theta}. Then $e_1 \cup e_2$ bounds a face of $M$.
\end{lemma}

\begin{proof}
Since $u \neq v$ and $e_1,e_2$ meet only at $u$ and $v$, the union
$\gamma = e_1 \cup e_2$ is a simple closed curve on the sphere. By
separation, its complement has exactly two components $D_1, D_2$, each an
open disc with boundary $\gamma$, and $\gamma$ carries no vertex other than
$u$ and $v$.

Cubicity leaves $u$ exactly one dart $p_a$ off $\gamma$, and $v$ exactly one
dart $q_b$. Each points into $D_1$ or into $D_2$.

\emph{They cannot point into different discs.} Say $p_a$ points into $D_1$
and $q_b$ into $D_2$, and let $a$ be the edge carrying $p_a$, with other end
$u'$. By Lemma~\ref{lem:no-loops}, $u' \neq u$. If $u' = v$ then $a$ is a
third $uv$ edge, saturating both vertices, and connectedness makes $M$ the
theta map --- excluded. So $u' \in D_1$.

Now every walk from $u'$ to $v$ must leave $D_1$, and it can only leave
through $\gamma$; but $\gamma$ carries no vertices besides $u$ and $v$, and
the only dart at either of them pointing into $D_1$ is $p_a$. So every such
walk traverses $a$, and $M - a$ separates $u'$ from $v$: the edge $a$ is a
bridge, contradicting bridgelessness.

\emph{Hence both point into the same disc}, say $D_1$. Then no dart at $u$
or at $v$ points into $D_2$. A vertex in the interior of $D_2$ would reach
the rest of $M$ only through $\gamma$, hence through such a dart, and there
is none; connectedness then forces the interior of $D_2$ to be free of
vertices, so also of edges. Thus $D_2$ is a face, bounded by $e_1$ and $e_2$
and nothing else.
\end{proof}

\subsection{The digon face, without separation}

The proof above spends a separation statement. It need not. The same
conclusion follows from two facts already machine-checked here --- the
bridge characterisation and Euler's \emph{inequality} --- and from the
observation that a digon in a cubic map is attached to the rest of the map
by exactly two edges.

\begin{theorem}[the digon face, combinatorially]\label{thm:digon-combinatorial}
\Sconditional{Theorem~\ref{thm:same-face-bridge}, Theorem~\ref{thm:euler-bound}}
Let $M$ be a connected bridgeless cubic map carrying a parallel pair
$e_1, e_2$ between distinct vertices $u, v$, and suppose $M$ is not the theta
map. If $M$ is spherical then
$\{e_1, e_2\}$ is a face of $M$.
\end{theorem}

\begin{proof}
Write $p_1, p_2, p_a$ for the darts at $u$ and $q_1, q_2, q_b$ for those at
$v$, with $\alpha p_i = q_i$; let $a$ be the edge carrying $p_a$, with second
dart $a'$, and $b$ the edge carrying $q_b$, with second dart $b'$. Since $u$
and $v$ have no further darts, every other dart lies in the exterior, and $M$
meets the exterior only along $a$ and $b$.

Relabelling $e_1 \leftrightarrow e_2$ if necessary, take
$\sigma_u = (p_1\,p_2\,p_a)$. The rotation at $v$ is one of two $3$-cycles.
If $\sigma_v = (q_1\,q_b\,q_2)$ then $\sigma q_2 = q_1$, and
\[
  \varphi(q_1) = \sigma p_1 = p_2, \qquad \varphi(p_2) = \sigma q_2 = q_1,
\]
so $\{q_1, p_2\}$ is a face of length two and we are done. Assume therefore
$\sigma_v = (q_1\,q_2\,q_b)$, and write $M^{\circ}$ for the map obtained by
replacing $\sigma_v$ with the other $3$-cycle. We show $M$ is not spherical.

Since $\varphi(x) = \sigma(\alpha x)$ and the two maps differ only in
$\sigma_v$, they differ only at darts $x$ with $\alpha x \in \{q_1,q_2,q_b\}$,
that is, at $p_1$, $p_2$ and $b'$. In particular the $\varphi$-orbit of a
dart of the exterior follows the same course in both maps until it reaches
$a'$ or $b'$.

Let $P_a$ be the segment from $\varphi(p_a) = \sigma a'$ up to its first
arrival at $a'$ or $b'$, and $P_b$ the segment from
$\varphi(q_b) = \sigma b'$. If $P_a$ arrived at $a'$ then $p_a$ and $a'$ ---
the two darts of $a$ --- would lie on one face, making $a$ a bridge by
Theorem~\ref{thm:same-face-bridge}; likewise an arrival of $P_b$ at $b'$
would make $b$ a bridge. Bridgelessness therefore forces $P_a$ to end at
$b'$ and $P_b$ at $a'$.

Now trace both maps. In $M$,
\[
  p_a \to P_a \to b' \to q_1 \to p_2 \to q_b \to P_b \to a' \to p_1 \to q_2
  \to p_a ,
\]
a single face. In $M^{\circ}$ the same darts fall into three:
\[
  (p_a,\, P_a,\, b',\, q_2), \qquad
  (p_1,\, q_b,\, P_b,\, a'), \qquad
  (q_1,\, p_2).
\]
Every face meeting neither $P_a$ nor $P_b$ is untouched, so
$F(M^{\circ}) = F(M) + 2$. The two maps share their underlying graph, hence
their vertex count, edge count and number of components. Euler's inequality
(Theorem~\ref{thm:euler-bound}) applied to $M^{\circ}$ gives
\[
  V + F(M) + 2 \;=\; V + F(M^{\circ}) \;\le\; E + 2 ,
\]
so $V + F(M) \le E$, and $M$ is not spherical.
\end{proof}

\begin{remark}[the last embedding obligation, discharged]
\Sprose
Theorem~\ref{thm:digon-combinatorial} uses no curve and no separation. Its
inputs are the bridge characterisation and Euler's inequality, both
machine-checked, together with a dart computation. Two features of the
situation carry it: a digon in a cubic map is attached by exactly two edges,
which is what makes the exterior enter only as two opaque segments $P_a$ and
$P_b$; and sphericity is the \emph{equality} case of an inequality that holds
unconditionally, so exhibiting a rival rotation with two more faces refutes
sphericity outright. The second is worth noting on its own --- an inequality
one cannot prove tight is still a tool for proving a given map is
\emph{not} tight.

An exhaustive check accompanies the theorem: over three cubic multigraphs
carrying a digon, all $16 + 64 + 256$ rotation systems were enumerated, and
no spherical bridgeless one failed to have the pair as a face.

Lemma~\ref{lem:digon-face} is retained for the picture it gives, but nothing
depends on it, and with it goes the last obligation in this development that
mentioned a plane.
\end{remark}

\begin{remark}[what this costs]
\Sprose
Nothing beyond what the route already owes.  The load-bearing proof is
Theorem~\ref{thm:digon-combinatorial}: bridge characterisation, Euler's
inequality, and a dart count.  Lemma~\ref{lem:digon-face} remains only as the
geometric picture and is not consumed.  The exclusion of the theta map is
harmless, since by Remark~\ref{rem:theta-colourable} it is Tait colourable and
so is never a least counterexample.
\end{remark}

\begin{lemma}[no parallel edges]\label{lem:no-parallel}
\Sprose
$M$ has no two vertices joined by more than one edge.
\end{lemma}

\begin{proof}
Suppose it does.  By cubicity there are either exactly two or exactly three
edges between the repeated pair of vertices.  In the latter case those two
vertices have no other incident darts, so connectedness makes $M$ the theta
map, which is Tait colourable by Remark~\ref{rem:theta-colourable}; this
contradicts that $M$ is a counterexample.  Hence there are exactly two, $M$ is
not theta, and we may form $M'$ as in Lemma~\ref{lem:digon-splice}. Then $M'$
lies in the same class and has two fewer vertices, so by minimality of $M$
it admits a Tait colouring $\chi'$.

Lift $\chi'$ to $M$. Give $a$ and $b$ the colour $\chi'(c)$, and give every
other edge of $M$ its $\chi'$-colour. It remains to colour the parallel
pair. At $u$ the edge $a$ already carries $\chi'(c)$, so assigning the two
parallel edges the two remaining colours makes the star at $u$ proper; the
same two colours make the star at $v$ proper, since $b$ also carries
$\chi'(c)$. Every other star is unchanged from $\chi'$ --- at $u'$ the edge
$c$ has been replaced by $a$ of the same colour, and likewise at $v'$.

So $M$ admits a Tait colouring, contradicting the assumption that it is a
counterexample.
\end{proof}

\subsection{The presentation}

\begin{theorem}[graph-backed presentation of a least counterexample]
\label{thm:graph-backed}
\Sprose
A least counterexample $M$ is the rotation system of a simple graph on its
own vertex set, and Tait colourings transfer between the two descriptions.
\end{theorem}

\begin{proof}
By Lemmas~\ref{lem:no-loops} and~\ref{lem:no-parallel} the incidence
relation of $M$ has no loop and no repeated pair, so it \emph{is} the
adjacency relation of a simple graph $G$ on the vertex set of $M$, and the
vertex rotations of $M$ are rotations of the darts of $G$. The edge sets are
in bijection under this identification, so a colouring of one is literally a
colouring of the other, and properness at a vertex is the same condition on
both sides.
\end{proof}

\begin{remark}[what the restriction bought]
\Sprose
The corresponding statement over the whole class is false by
Example~\ref{ex:theta}. What minimality supplies is precisely
Lemma~\ref{lem:no-parallel}, and it supplies it in the only way available:
by producing a smaller member of the class and colouring it. This is worth
emphasising because the formal statement must carry the minimality data as a
hypothesis --- a version quantified over all bridgeless spherical cubic maps
is not merely unproved but refuted, and a development asserting it would be
unsound rather than incomplete.
\end{remark}

\begin{remark}[status summary for the classical boundary]
\Sprose
With this section the ordinary mathematics of the classical bridge is closed.
Two-sidedness and both directions of the bridge characterisation are
machine-checked.  Theorem~\ref{thm:digon-combinatorial} removes the only
geometric premise from the digon splice, so minimality eliminates loops and
parallel edges and yields Theorem~\ref{thm:graph-backed} without a separation
axiom.

The companion reduction in Section~\ref{sec:tait-reduction} likewise uses no
maximal-plane-graph argument: stellar subdivision is an explicit dart
substitution, duality is the replacement of $\sigma$ by $\sigma\alpha$, and
looplessness of the subdivision makes its dual bridgeless.  The classical
maximal triangulation lemma and the geometric digon lemma remain in the text
as recognizable comparison routes; their \textsc{open}/\textsc{conditional}
markers describe those statements themselves, not dependencies of the
controlling proof.  What remains here is Lean transport and construction, not
an ordinary-mathematics gap.
\end{remark}
